\section{Companion lab: a sprite that crosses the frame boundary}
\label{revision-lab:clipped-sprite}

The row renderer above receives a destination where every output pixel fits.
A movable sprite has another problem: its origin can lie outside the frame while
some of its pixels remain visible. Moving it two columns left should hide those
two source columns and preserve the rest. Moving it completely beyond an edge
should change no destination pixel. Neither case permits writing outside the
frame and hoping that the display will hide the mistake.

This companion develops original project \texttt{016-clipped-sprite-blitter}
through the native \texttt{016-clipped-regions} unit. Open its assignment and
begin with \texttt{worked/interval.asm}, a complete horizontal clipping experiment.
The second experiment, \texttt{worked/extreme.asm}, changes only the placement to
the largest signed word. Compare their intermediate endpoints before implementing
the image copy.

Your implementation belongs in \texttt{independent/blit.asm}. Replace the failing
\texttt{blit\_sprite} body while retaining the complete caller, allocations,
descriptors and public observation labels. A \emph{blit} is a rectangular image
copy; this version also clips the copy and treats a selected source word as
transparent. First check the empty-intersection path, then complete the overlapping
copy and check the caller. Choose your private labels, register allocation and
equivalent algorithm. The original archived project remains available under its
package path.

Use the chapter's row-major addressing, the earlier signed comparisons and the
function calling convention. The new reasoning is to intersect rectangles,
translate the surviving destination origin back into source coordinates, and
keep logical width separate from physical row stride. The interval experiment
below is complete; the full image copy is your independent implementation.

\subsection*{Read the image descriptors before forming addresses}

The function receives a framebuffer descriptor in \texttt{\$a0}, a sprite
descriptor in \texttt{\$a1}, and placement coordinates $x,y$ in
\texttt{\$a2}/\texttt{\$a3}. The coordinates are signed 32-bit values, already
sign extended to the full register width. They locate the sprite's upper-left
cell in the frame's coordinate system. Both coordinate systems increase rightward
and downward.

\begin{table}[htbp]
  \centering
  \small
  \begin{tabular}{rlll}
    \toprule
    Byte offset & Size       & Frame descriptor      & Sprite descriptor     \\
    \midrule
    0           & doubleword & pixel-storage address & pixel-storage address \\
    8           & word       & logical width         & logical width         \\
    12          & word       & logical height        & logical height        \\
    16          & word       & row stride in words   & row stride in words   \\
    20          & word       & not a frame field     & transparency key      \\
    \bottomrule
  \end{tabular}
  \caption{Descriptors are eight-byte aligned; their pixel pointers are four-byte aligned. Only the sprite descriptor supplies a key. Words following the frame fields may belong to a guard, not to another readable field.}
  \label{tab:revision-clipped-descriptors}
\end{table}


A descriptor's width says how many pixels belong to a row. Its stride says how
many words separate the starts of adjacent rows, including any padding. For a
pixel at local column $c$ and row $r$, its address is
\[
  base+4(r\cdot stride+c).
\]
The multiplier four converts words to bytes. It does not convert width into
stride: a width-five sprite with stride seven has two padding words after every
logical row. Those words are storage, but they are not source pixels.

Dimensions range from zero through eight, and stride is at least width and at
most sixteen. All logical cells fit within the supplied 128-word allocation;
the stored pointer can name an interior word rather than the allocation's first
label. The two pixel allocations and descriptors are disjoint. You can assume
valid descriptors; repairing corrupt dimensions and supporting overlapping copies
are separate problems.

Load the address field as a doubleword and add offsets with full-width arithmetic.
Load a pixel as a word and copy its complete 32-bit pattern. A pointer and a pixel
have different widths even when both happen to be held in general-purpose
registers. Derive each address from the passed descriptor, rather than substituting
the caller's default \texttt{frame\_a} or \texttt{sprite\_a} label.

\subsection*{Intersect half-open intervals in frame coordinates}

Write the frame's horizontal interval as $[0,F_w)$ and the placed sprite's as
$[x,x+S_w)$. A \emph{half-open interval} includes its left endpoint and excludes
its right endpoint. Its length is the difference between the endpoints, so adjacent
intervals can meet without sharing a pixel. The surviving destination interval is
\[
  L=\max(0,x),\qquad R=\min(F_w,x+S_w).
\]
It is nonempty only when $L<R$. Repeat for rows:
\[
  T=\max(0,y),\qquad B=\min(F_h,y+S_h),
\]
requiring $T<B$. If either comparison fails, return zero visible and zero opaque
pixels without a destination store. A zero width or height is an empty rectangle.
A sprite beginning exactly at $x=F_w$ shares no columns with the frame.

Keep these endpoint calculations at full width. Although each image is small,
$x$ and $y$ can be extreme signed words. Their sums with a dimension and the
subsequent differences need not fit a signed word. With $x=2147483647$ and width
five, the mathematical right endpoint is 2147483652. A word addition followed by
a word subtraction can turn a rejected interval into a plausible small positive
extent. Prove the intersection empty before deriving pixel addresses.

The complete \texttt{worked/interval.asm} experiment applies the horizontal rule.
It receives $x$ in \texttt{\$a0}, source width in \texttt{\$a1}, and frame width
in \texttt{\$a2}; it returns $L,R$ in \texttt{\$v0}/\texttt{\$v1}, or two
zeros for an empty intersection. This is a helper interface for the experiment,
not the blitter's pixel-count return interface.

\noindent\begin{minipage}{\linewidth}
  \MARSCodeLabel{Select the smaller right endpoint using full-width arithmetic}
  \begin{minted}{mars-idea}
interval_left_ready:
    daddu $t1, $a0, $a1
    slt $t2, $a2, $t1
    beqzc $t2, interval_right_ready
    nop
    move $t1, $a2
\end{minted}
\end{minipage}

Here \texttt{\$t1} begins as the placed source's right endpoint. If frame width
is smaller, the move selects it; otherwise the branch preserves the endpoint.
The complete helper also selects the larger left endpoint and rejects $L\geq R$.
Run its public \texttt{worked/interval.kts} check to cover 2,268 argument combinations and the complete default caller. Predict the
supplied $x=-2$, source width five and frame width eight: the result is $[0,3)$.
The guided Debug sequence observes the selected left endpoint, selected right
endpoint and returned interval. In \texttt{worked/extreme.asm}, the placement is
2147483647. Pause immediately after \texttt{daddu}: the source end is 2147483652,
hexadecimal \texttt{0000000080000004}. It has crossed the signed-word maximum
without becoming a negative full-register value. The selected right endpoint then
becomes frame width eight, while the left remains 2147483647. That ordering rejects
the interval and returns two zeros. No pixel address is needed.

Also test a source ending at zero and one beginning at eight. Explain the
comparison that makes each interval empty. Both complete experiments exit
silently after recording their endpoints; their console transcript alone does
not establish those values.

\subsection*{Translate the surviving origin back into the source}

The first visible destination cell is $(L,T)$. It corresponds to source cell
\[
  s_x=L-x,\qquad s_y=T-y.
\]
The visible width is $w=R-L$ and height is $h=B-T$. Clipping the destination does
not move source column zero to $L$: it discards the columns and rows before
$(s_x,s_y)$. Those source offsets are the reason the visible fragment retains the
sprite's original shape.

The default sprite is five by four with stride seven, placed at $(-2,1)$ in an
eight-by-six frame with stride sixteen. The intersection is $[0,3)\times[1,5)$,
so $(s_x,s_y)=(2,0)$ and twelve source cells are visible.
Figure~\ref{fig:revision-clipped-geometry} connects those coordinates.

\begin{figure}[htbp]
\centering
\begin{tikzpicture}[x=7mm,y=-7mm,alt={A five by four sprite is placed two columns left and one row down from an eight by six frame. Two source columns lie outside the frame. The visible destination rectangle covers columns zero through two and rows one through four, corresponding to source columns two through four and rows zero through three. A solid outline marks the frame, a dashed outline marks the complete placed sprite, and shading marks the intersection.}]
\fill[black!15] (0,1) rectangle (3,5);
\foreach \x in {-2,-1,0,1,2,3} {\draw[black!35] (\x,1)--(\x,5);}
\foreach \y in {1,...,5} {\draw[black!35] (-2,\y)--(3,\y);}
\foreach \x in {0,...,8} {\draw[black!35] (\x,0)--(\x,6);}
\foreach \y in {0,...,6} {\draw[black!35] (0,\y)--(8,\y);}
\draw[very thick] (0,0) rectangle (8,6);
\draw[thick,dashed] (-2,1) rectangle (3,5);
\draw[thick] (0,1) rectangle (3,5);
\foreach \x in {0,...,7} {\node[font=\small,anchor=south] at (\x+.5,-.12) {\x};}
\foreach \y in {0,...,5} {\node[font=\small,anchor=west] at (8.15,\y+.5) {\y};}
\node[anchor=north,font=\small] at (1.5,6.35) {Visible destination columns 0--2};
\node[anchor=east,align=right,font=\small] at (-2.2,3) {Hidden source\\columns 0--1};
\node[anchor=west,align=left,font=\small,text width=42mm] at (9,2.6) {$L=0,\ R=3$\\$T=1,\ B=5$\\Source origin: $(2,0)$};
\end{tikzpicture}
\caption{Logical geometry for placement $(-2,1)$. The dashed sprite outline crosses the solid frame boundary. Physical row padding is omitted from this coordinate diagram; it remains present in storage.}
\label{fig:revision-clipped-geometry}
\end{figure}


The first source address is $source+4(0\cdot7+2)=source+8$. The first destination
address is $frame+4(1\cdot16+0)=frame+64$. The next visible source row begins
28 bytes after the first source row; the next destination row begins 64 bytes
after its predecessor. Neither increment is the twelve-byte visible row length.

Keep row starts separate from moving pixel pointers. One organization stores each
row's source and destination starts, copies $w$ positions, then advances the row
starts by their respective byte strides. Another lets the inner pointers advance
and then adds $4(stride-w)$ at the end of each row. Using a full stride after an
already-advanced inner pointer skips the copied width a second time.

Use a loop invariant to distinguish these stages. At row $j$ and column $i$ within
the intersection, the next selected source cell is $(s_x+i,s_y+j)$ and the next
destination cell is $(L+i,T+j)$. Before the first pixel, $i=0$; after exactly $w$
positions the row is complete. A transparent pixel still occupies one position
and advances both pointers.

You may instead visit every logical source cell and test its translated destination
coordinates before copying. The contract accepts that organization too. Source
reads may inspect logical source cells that are later clipped away, but never
source padding or guards. Destination reads, if used, must remain within logical
frame cells. There is no prescribed read count, register allocation or traversal
order. The precomputed intersection is a way to make the bounds explicit, not a
required source-code pattern.

\subsection*{Separate visibility, transparency and changed colors}

Return the number of visible opaque source pixels in \texttt{\$v0}, and all
visible source pixels in \texttt{\$v1}. The latter is $wh$, including transparent
cells. A fully invisible or empty image returns $(0,0)$. These counts describe
selected source cells, not the number of destination colors that changed.

The transparency key is the complete 32-bit word in the sprite descriptor. A
source word equal to that key skips its destination store. Every other word is
opaque and must be copied unchanged. Zero is opaque unless zero is the selected
key. The high byte is part of the comparison; this lab does not interpret it as
an alpha channel.

In the supplied sprite, key \texttt{0x00FF00FF} is transparent, while
\texttt{0x01FF00FF} is opaque. The Bitmap Display uses the low RGB bytes, so those
two words can look the same when displayed directly. Their different copy behavior
comes from the word comparison. Inspect Memory when the display cannot distinguish
the relevant bits.

\begin{table}[htbp]
  \centering
  \small
  \begin{tabular}{rlll}
    \toprule
    Source row & Source column 2   & Source column 3   & Source column 4   \\
    \midrule
    0          & key: no store     & \texttt{0000FF00} & key: no store     \\
    1          & \texttt{00FFFF00} & \texttt{01FF00FF} & \texttt{00FF0000} \\
    2          & \texttt{0000FFFF} & key: no store     & \texttt{0000FF00} \\
    3          & \texttt{000000FF} & \texttt{00FFFF00} & key: no store     \\
    \bottomrule
  \end{tabular}
  \caption{The default intersection selects these twelve source words, written in hexadecimal. Source rows zero through three map to destination rows one through four. Four keys preserve the previous destination words; eight opaque cells are copied.}
  \label{tab:revision-clipped-cells}
\end{table}


For the default placement, the twelve visible cells include four keys and eight
opaque words, so return $(8,12)$. The opaque word in the second row of the table
has the same low RGB bytes as the key and must still be stored. With placement
$(1,1)$, all twenty source cells fit and thirteen are opaque. One of those cells
contains zero and must draw black. Use these two placements to distinguish a
complete-word comparison from an RGB-only comparison or a zero-terminated rule.

An opaque source pixel counts even if the destination already has exactly that
word. A routine that counts only changed colors reports the wrong interface.
Conversely, a transparent destination must receive no store, including a store
of its old value or a write later restored. A final matching picture cannot
establish that rule: the grading checks the store history as well as the frame.
Preserve clipped cells, row padding, guards, descriptors and the source throughout.

\subsection*{Observe the copy and diagnose an incorrect span}

At \texttt{blit\_call}, the caller has initialized the frame and both descriptor
pointers but has not entered your function. At \texttt{blit\_ready}, the function
has returned; inspect its result registers before the caller publishes
\texttt{copied} and \texttt{visible}. After normal exit, those words retain the
counts, while the exit service has replaced \texttt{\$v0}.

The revised source requests a sixteen-by-eight Bitmap Display beginning at
\texttt{frame\_a} in the Default memory layout. The logical frame is only the
first eight columns of its first six rows. The remaining displayed cells are
allocation padding, prefilled gray; the active frame is prefilled dark blue.
Display row width matches stride sixteen so padding does not masquerade as a
shorter next row. If you change the memory layout or data origin, select the new
frame base in the display too. Display configuration does not alter a descriptor.

Figure~\ref{fig:revision-clipped-observations} compares the default partial copy
with a complete placement. Native outlines identify the logical frame independently
of the background colors; the surrounding storage is padding.

\begin{figure}[htbp]
\centering
\begin{tikzpicture}[alt={Two actual light-theme Bitmap Display observations of a sixteen by eight storage view. A double outline marks the logical eight by six frame in its upper-left region; the gray right half and bottom two rows are padding. First panel shows placement minus two comma one, with eight opaque copies among twelve visible cells. Second panel shows placement one comma one, with thirteen opaque copies among twenty visible cells. Native marker one identifies the copied word 01FF00FF in both panels; marker two identifies the opaque zero word in the second panel.}]
\node[anchor=north west,inner sep=0,font=\small\bfseries] (atitle) at (0,0) {1. Placement $(-2,1)$: left columns clipped};
\node[anchor=north west,inner sep=0] (a) at ([yshift=-2mm]atitle.south west) {\MARSRevisionCrop{revision/captures/clipped/clipped-left.png}{0.82\linewidth}{0.032031250}{0.155555556}{0.032031250}{0.178888889}};
\node[anchor=north west,inner sep=0,font=\small\bfseries] (btitle) at ([yshift=-4mm]a.south west) {2. Placement $(1,1)$: complete source visible};
\node[anchor=north west,inner sep=0] (b) at ([yshift=-2mm]btitle.south west) {\MARSRevisionCrop{revision/captures/clipped/clipped-full.png}{0.82\linewidth}{0.032031250}{0.155555556}{0.032031250}{0.178888889}};
\foreach \panel in {a,b} {
\begin{scope}[shift={(\panel.south west)},x={($(\panel.south east)-(\panel.south west)$)},y={($(\panel.north west)-(\panel.south west)$)}]
\draw[black,line width=1.8pt] (0,.25) rectangle (.5,1);
\draw[white,line width=.7pt] (0,.25) rectangle (.5,1);
\end{scope}}
\foreach \panel/\x/\y/\n in {a/.09375/.6875/1,b/.28125/.6875/1,b/.15625/.6875/2} {
\begin{scope}[shift={(\panel.south west)},x={($(\panel.south east)-(\panel.south west)$)},y={($(\panel.north west)-(\panel.south west)$)}]
\node[font=\scriptsize\bfseries,fill=white,text=black,draw=black,circle,inner sep=1pt] at (\x,\y) {\n};
\end{scope}}
\end{tikzpicture}
\caption{Actual completed copies. The native publication outline separates logical frame cells from allocation padding. Marker 1 identifies an opaque word whose low RGB bytes match the key; marker 2 identifies an opaque zero that draws black. Key cells retain the previous background. All 128 native display cells were checked for each observation.}
\label{fig:revision-clipped-observations}
\end{figure}


Begin by drawing the default placement on paper, recording the two first addresses
and the twelve visible source words. Run your implementation and compare its
frame and counts. Then change placement to $(6,-2)$: source rows two and three,
columns zero and one survive, mapping to destination columns six and seven in
rows zero and one. There are four visible cells, two opaque; return $(2,4)$.

In disposable copies, replace one row stride with logical width and locate the
first address that diverges. Next retain the correct strides but set the clipped
source origin to zero. A fragment appears in the expected place but contains the
wrong source columns. For transparency, compare only the low 24 bits and observe
the missing opaque magenta cell. Diagnose each fault from the first wrong source
selection, address or store, rather than adjusting a counter until a check passes.

The function must preserve full saved registers \texttt{\$s0} through
\texttt{\$s7}, \texttt{\$fp} and \texttt{\$sp}, return through the supplied
\texttt{\$ra}, and produce no console output. A restored private stack frame is
permitted. If you call the interval helper, preserve the blitter's return link and
any live state that its caller-clobbered temporaries may overwrite. Inlining the
same comparisons is also valid.

Use the native \emph{Check this program} control first for the 474
empty-intersection cases. You can finish safe geometric rejection while leaving
the nonempty copy path unfinished. Empty dimensions, exact edge contact,
fully clipped placements and extreme coordinates all belong to this first domain.
They return $(0,0)$ without a destination store; a frame being allocated does not
make it writable for an invisible sprite.

Then complete the copy and check the 107 overlap cases in the same source. Together
the two assessments retain all 581 original kernel trials. Keep the geometry,
source translation and copy invariant in view when interpreting a failure. A
correct visible count with incorrect pixels suggests a different problem from a
wrong extent or a forbidden transparent store.

The ten caller cases observe initialized storage at \texttt{blit\_call}, then
returned registers and frame at \texttt{blit\_ready}, then stored counts after
normal exit. The program continues between those boundaries without resetting
registers. A wrong return value corrected only before publication still violates
the function interface. The final observation interval also protects caller data
after return.

For the guided Debug activity, place breakpoints at \texttt{blit\_call} on
\texttt{balc}, \texttt{blit\_ready} on \texttt{nop}, and \texttt{done} on
\texttt{li}. Inspect those stops in order in the same session. Before the call,
placement registers hold $(-2,1)$ and the frame has its initial colors. At return,
$(v_0,v_1)=(8,12)$ while the caller has not yet published the counts. Before the
exit selector, the result registers still hold those values and the memory words
\texttt{copied} and \texttt{visible} contain eight and twelve. Stop Debug before
the final ordinary Run.

The public scripts also run from the downloaded unit's directory using JDK 25.
The caller script uses the same observation engine as the native assessment and
requires the MARS edition supporting console format 8.

\noindent\begin{minipage}{\linewidth}
  \begin{verbatim}
java -jar /path/to/cli.jar grade \
  independent/kernel.kts independent/blit.asm
java -jar /path/to/cli.jar grade \
  independent/caller.kts independent/blit.asm
\end{verbatim}
\end{minipage}

Run normally and expect an empty console. Compare image, memory and stored counts.
Keep an empty intersection, a partial overlap and an all-transparent example.
Explain the last example's positive visible count, zero opaque count and protected
padding.
